\section{Limits, Infinitesimals, Continuity, and the Derivative}

In arithmetic, we enlarged the number system whenever an equation could not be solved.

In calculus, a different kind of difficulty appears.

Sometimes we want to understand a value that is not reached directly, but is approached.

For example, consider the numbers

\[
2,\quad 1.5,\quad 1.333\ldots,\quad 1.25,\quad 1.2,\quad \ldots
\]

They seem to move closer and closer to (1), but no term in this list is exactly (1).

This raises a new question.

\begin{quote}
How can a variable have a final tendency even if it never reaches that value?
\end{quote}

To study this question, consider the sequence

\[
x_n = 1+\frac{1}{n}.
\]

When (n=1),

\[
x_1=2.
\]

When (n=2),

\[
x_2=\frac{3}{2}.
\]

When (n=10),

\[
x_{10}=1.1.
\]

When (n=1000),

\[
x_{1000}=1.001.
\]

The values are not equal to (1), but they become closer and closer to (1).

So the important object is not equality.

The important object is distance.

The distance from $$(x_n)$$ to (1) is

\[
|x_n-1|.
\]

Computing,

\[
|x_n-1| = \left|1+\frac{1}{n}-1\right| \frac{1}{n}.
\]

Now we ask:

\begin{quote}
Can this distance be made smaller than any positive number chosen in advance?
\end{quote}

Suppose someone gives us a small number $(\varepsilon>0)$.

We want

\[
|x_n-1|<\varepsilon.
\]

Since

\[
|x_n-1|=\frac{1}{n},
\]

we need

\[
\frac{1}{n}<\varepsilon.
\]

This is true whenever

\[
n>\frac{1}{\varepsilon}.
\]

Therefore, no matter how small $(\varepsilon)$ is, we can choose (n) large enough so that all later values of $(x_n)$ are within distance $(\varepsilon)$ of (1).

This phenomenon deserves a name.

\begin{definition}
A number (a) is called the \emph{limit} of a sequence $(x_n)$ if, for every
$(\varepsilon>0)$, there exists an integer (N) such that whenever (n$>$N),

\[
|x_n-a|<\varepsilon.
\]

In this case we write

\[
\lim_{n\to\infty}x_n=a.
\]
\end{definition}

Thus,

\[
\lim_{n\to\infty}\left(1+\frac{1}{n}\right)=1.
\]

\subsection{Problems}

\subsubsection*{Problem 1}

Consider

\[
x_n=\frac{1}{n}.
\]

What number does $(x_n)$ approach as $(n\to\infty)$?

\subsubsection*{Problem 2}

Show that

\[
\lim_{n\to\infty}\frac{1}{n}=0.
\]

That is, given $(\varepsilon>0)$, find a condition on (n) which guarantees

\[
\left|\frac{1}{n}-0\right|<\varepsilon.
\]

\subsubsection*{Problem 3}

Consider

\[
x_n=3+\frac{2}{n}.
\]

What is the limit?

Prove your answer using the definition.

\section{The Limit of a Function}

Sequences are not the only objects that approach values.

A function can also approach a value.

Consider

\[
f(x)=\frac{x^2-1}{x-1}.
\]

At first, something strange happens when (x=1).

If we substitute (x=1), we get

\[
f(1)=\frac{1^2-1}{1-1}=\frac{0}{0},
\]

which is not defined.

So the function has no value at (x=1).

But what happens near (x=1)?

Since

\[
x^2-1=(x-1)(x+1),
\]

we have, for $(x\neq 1)$,

\[
f(x)=\frac{(x-1)(x+1)}{x-1}=x+1.
\]

Therefore, when (x) is close to (1), the values of (f(x)) are close to

\[
1+1=2.
\]

For example,

\[
f(0.9)=1.9,
\]

\[
f(0.99)=1.99,
\]

\[
f(1.001)=2.001.
\]

Even though (f(1)) is not defined, the nearby values of (f(x)) clearly approach (2).

This creates a new question.

\begin{quote}
Can a function approach a value at a point even if the function is not defined at that point?
\end{quote}

The answer is yes.

This is why we need the limit of a function.

\begin{definition}
Let (f(x)) be a function. We say that

\[
\lim_{x\to a}f(x)=b
\]

if, for every $(\varepsilon>0)$, there exists a $(\delta>0)$ such that whenever

\[
0<|x-a|<\delta,
\]

we have

\[
|f(x)-b|<\varepsilon.
\]
\end{definition}

The condition

\[
0<|x-a|
\]

means that (x) is allowed to approach (a), but (x) is not required to equal (a).

This is essential.

A limit is about nearby behavior.

It is not necessarily about the value at the point.

\subsection{Example}

Let

\[
f(x)=\frac{x^2-1}{x-1}.
\]

For $(x\neq 1)$,

\[
f(x)=x+1.
\]

Therefore,

\[
\lim_{x\to 1}\frac{x^2-1}{x-1} = \lim_{x\to 1}(x+1) 2.
\]

So the function does not have a value at (x=1), but it has a limit there.

\subsection{Problems}

\subsubsection*{Problem 4}

Compute

\[
\lim_{x\to 2}\frac{x^2-4}{x-2}.
\]

First simplify the expression for $(x\neq 2)$.

\subsubsection*{Problem 5}

Consider

\[
f(x)=
\begin{cases}
x+1, & x\neq 1,\\
10, & x=1.
\end{cases}
\]

What is

\[
\lim_{x\to 1}f(x)?
\]

What is (f(1))?

Are they equal?

\section{Infinitesimals}

When (x) approaches (a), the difference

\[
x-a
\]

approaches (0).

When (f(x)) approaches (b), the difference

\[
f(x)-b
\]

also approaches (0).

So many limit questions reduce to studying quantities that become smaller and smaller.

This suggests another name.

\begin{definition}
A function $(\alpha(x))$ is called an $\emph{infinitesimal}$ as $(x\to a)$ if

\[
\lim_{x\to a}\alpha(x)=0.
\]

Similarly, $(\alpha(x))$ is called an infinitesimal as $(x\to\infty)$ if

\[
\lim_{x\to\infty}\alpha(x)=0.
\]
\end{definition}

For example,

\[
\frac{1}{x}
\]

is infinitesimal as $(x\to\infty)$, because

\[
\lim_{x\to\infty}\frac{1}{x}=0.
\]

Also,

\[
x-a
\]

is infinitesimal as $(x\to a)$, because

\[
\lim_{x\to a}(x-a)=0.
\]

\subsection{A Limit Written as a Constant Plus an Error}

Suppose a variable (y) approaches (b).

Then the difference

\[
y-b
\]

approaches (0).

So (y-b) is infinitesimal.

Let

\[
\alpha=y-b.
\]

Then

\[
y=b+\alpha.
\]

This means that a quantity approaches (b) exactly when it can be written as

\[
\text{constant}+\text{infinitesimal}.
\]

\subsubsection*{Theorem 1}

A variable (y) satisfies

\[
\lim y=b
\]

if and only if it can be written in the form

\[
y=b+\alpha,
\]

where $(\alpha)$ is infinitesimal.

\paragraph{Proof.}

If

\[
y=b+\alpha
\]

and $(\alpha\to 0)$, then

\[
y-b=\alpha.
\]

Therefore,

\[
|y-b|=|\alpha|.
\]

Since $(\alpha)$ is infinitesimal, $(|\alpha|)$ can be made smaller than any $(\varepsilon>0)$. Hence

\[
y\to b.
\]

Conversely, if $(y\to b)$, define

\[
\alpha=y-b.
\]

Then

\[
\alpha\to 0,
\]

so $(\alpha)$ is infinitesimal, and

\[
y=b+\alpha.
\]

\subsection{Algebra of Infinitesimals}

Once infinitesimals appear, we must ask whether they behave well under ordinary operations.

\begin{quote}
If two quantities become negligible, does their sum also become negligible?
\end{quote}

\subsubsection*{Theorem 2}

The sum of finitely many infinitesimals is infinitesimal.

\paragraph{Proof.}

Let

\[
u(x)=\alpha(x)+\beta(x),
\]

where both $(\alpha)$ and $(\beta)$ are infinitesimal.

Given $(\varepsilon>0)$, choose (x) close enough to the limiting value so that

\[
|\alpha(x)|<\frac{\varepsilon}{2}
\]

and

\[
|\beta(x)|<\frac{\varepsilon}{2}.
\]

Then

\[
|u(x)| =|\alpha(x)+\beta(x)| \leq |\alpha(x)|+|\beta(x)| < \frac{\varepsilon}{2}+\frac{\varepsilon}{2} = \varepsilon.
\]

Thus (u(x)) is infinitesimal.

\subsubsection*{Theorem 3}

The product of an infinitesimal and a bounded function is infinitesimal.

\paragraph{Proof.}

Let $(\alpha(x))$ be infinitesimal and let (z(x)) be bounded.

Since (z(x)) is bounded, there exists (M>0) such that

\[
|z(x)|\leq M.
\]

Given $(\varepsilon>0)$, choose (x) close enough so that

\[
|\alpha(x)|<\frac{\varepsilon}{M}.
\]

Then

\[
|z(x)\alpha(x)| = |z(x)|,|\alpha(x)| \leq M|\alpha(x)| < M\cdot \frac{\varepsilon}{M} = \varepsilon.
\]

Therefore $(z(x)\alpha(x))$ is infinitesimal.

\subsubsection*{Theorem 4}

If $(\alpha(x))$ is infinitesimal and $(\alpha(x)\neq 0)$, then

\[
\left|\frac{1}{\alpha(x)}\right|
\]

becomes arbitrarily large.

\paragraph{Proof.}

Let (M>0). Since $(\alpha(x)\to 0)$, we can choose (x) close enough so that

\[
|\alpha(x)|<\frac{1}{M}.
\]

Then

\[
\left|\frac{1}{\alpha(x)}\right| = \frac{1}{|\alpha(x)|} > M.
\]

So the reciprocal of a nonzero infinitesimal becomes arbitrarily large.

\section{Continuity}

We have seen that a function may have a limit at a point even if it is not defined there.

We have also seen that a function may be defined at a point but have a value different from the nearby behavior.

For example,

\[
f(x)=
\begin{cases}
x+1, & x\neq 1,\\
10, & x=1.
\end{cases}
\]

Then

\[
\lim_{x\to 1}f(x)=2,
\]

but

\[
f(1)=10.
\]

So the function has a break at (x=1).

This raises another question.

\begin{quote}
When does the value of a function agree with the value approached by nearby points?
\end{quote}

This is the idea of continuity.

\begin{definition}
A function (f(x)) is called $\emph{continuous}$ at $(x_0)$ if

\[
\lim_{x\to x_0}f(x)=f(x_0).
\]
\end{definition}

In words:

\begin{quote}
A function is continuous at $(x_0)$ when the value it approaches is the value it actually has.
\end{quote}

\subsection{Continuity Written with Increments}

Let

\[
x=x_0+\Delta x.
\]

Then

\[
\Delta x=x-x_0.
\]

As $(x\to x_0)$, we have

\[
\Delta x\to 0.
\]

The change in the function is

\[
\Delta y=f(x_0+\Delta x)-f(x_0).
\]

If (f) is continuous at $(x_0)$, then as $(\Delta x\to 0)$,

\[
\Delta y\to 0.
\]

Therefore, continuity can also be written as

\[
\lim_{\Delta x\to 0} \bigl(f(x_0+\Delta x)-f(x_0)\bigr) = 0.
\]

This says:

\begin{quote}
A small change in (x) produces a small change in (y).
\end{quote}

\subsection{Direct Substitution}

If (f) is continuous at $(x_0)$, then

\[
\lim_{x\to x_0}f(x)=f(x_0).
\]

Since

\[
x_0=\lim_{x\to x_0}x,
\]

we may write informally

\[
\lim_{x\to x_0}f(x) = f\left(\lim_{x\to x_0}x\right).
\]

This is why continuous functions allow limits to be computed by direct substitution.

\subsection{Problems}

\subsubsection*{Problem 6}

Let

\[
f(x)=x^2.
\]

Show that (f) is continuous at $(x_0=3)$.

That is, show that

\[
\lim_{x\to 3}x^2=9.
\]

\subsubsection*{Problem 7}

Let

\[
f(x)=
\begin{cases}
x^2, & x\neq 2,\\
7, & x=2.
\end{cases}
\]

Find

\[
\lim_{x\to 2}f(x).
\]

Find (f(2)).

Is (f) continuous at (2)?

\section{Continuity on an Interval}

Continuity at one point is a local property.

But sometimes we want a function to behave without breaks over a whole interval.

\begin{definition}
A function (f(x)) is continuous on an interval ((a,b)) if it is continuous at every point of that interval.
\end{definition}

If the interval is closed,

\[
[a,b]
\],

then (f) must be continuous at every interior point, continuous from the right at (a), and continuous from the left at (b).

\section{Properties of Continuous Functions}

Once continuity is understood as ``no breaks,'' we can ask what such functions are forced to do.

For example, suppose a function is continuous on a closed interval

\[
[a,b].
\]

Can it grow forever inside this interval?

No.

The interval is finite, and the function has no breaks that allow it to escape.

This leads to an important theorem.

\subsubsection*{Theorem 5 (Extreme Value Theorem)}

If (f(x)) is continuous on a closed interval ([a,b]), then there exist points ($x_1,x_2 \in$ [a,b]) such that

\[
f(x_1)\geq f(x)
\]

and

\[
f(x_2)\leq f(x)
\]

for every $(x\in [a,b])$.

In other words, a continuous function on a closed interval attains a maximum and a minimum.

\section{A Problem Before the Derivative: Fermat's Segment Problem}

Before defining the derivative, let us begin with a problem.

Suppose we have a segment (AB) of fixed length (L).

Choose a point (C) between (A) and (B).

Let

\[
AC=x.
\]

Then

\[
CB=L-x.
\]

We want to choose (C) so that the product

\[
AC\cdot CB
\]

is as large as possible.

In symbols, we want to maximize

\[
P(x)=x(L-x).
\]

\subsection{First Experiments}

If (C) is very close to (A), then (AC) is small.

So

\[
AC\cdot CB
\]

is small.

If (C) is very close to (B), then (CB) is small.

Again,

\[
AC\cdot CB
\]

is small.

So the maximum should happen somewhere in the middle.

Let us test some values.

If

\[
x=0,
\]

then

\[
P(0)=0(L-0)=0.
\]

If

\[
x=\frac L4,
\]

then

\[
P\left(\frac L4\right) = \frac L4\left(L-\frac L4\right) \frac L4\cdot \frac{3L}{4}\frac{3L^2}{16}.
\]

If

\[
x=\frac L2,
\]

then

\[
P\left(\frac L2\right) = \frac L2\left(L-\frac L2\right) \frac L2\cdot \frac L2 \frac{L^2}{4}.
\]

The middle seems better.

But experiments are not proof.

We need a method.

\subsection{Moving the Point Slightly}

Suppose the point (C) is already chosen.

Now move it a very small distance (h) to the right.

Then the new value of (AC) is

\[
x+h.
\]

The new value of (CB) is

\[
L-x-h.
\]

So the new product is

\[
P(x+h)=(x+h)(L-x-h).
\]

The old product was

\[
P(x)=x(L-x).
\]

The question becomes:

\begin{quote}
Did the product increase or decrease after this small movement?
\end{quote}

To answer this, compute the difference

\[
P(x+h)-P(x).
\]

We have

\[
P(x+h)-P(x) = (x+h)(L-x-h)-x(L-x).
\]

Expand the first product:

\[
(x+h)(L-x-h) = x(L-x-h)+h(L-x-h).
\]

So

\[
(x+h)(L-x-h) = x(L-x)-xh+h(L-x)-h^2.
\]

Therefore,

\[
P(x+h)-P(x) = x(L-x)-xh+h(L-x)-h^2-x(L-x).
\]

Canceling (x(L-x)), we get

\[
P(x+h)-P(x) = -xh+h(L-x)-h^2.
\]

Thus,

\[
P(x+h)-P(x) = h(L-2x)-h^2.
\]

\subsection{The Important Part of the Change}

The change in the product is

\[
P(x+h)-P(x) = h(L-2x)-h^2.
\]

Now notice something important.

If (h) is very small, then $(h^2)$ is much smaller than (h).

For example, if

\[
h=0.01,
\]

then

\[
h^2=0.0001.
\]

So the term

\[
h(L-2x)
\]

is the main part of the change.

The term

\[
-h^2
\]

is smaller.

This suggests that the expression

\[
L-2x
\]

controls whether the product first wants to increase or decrease.

\subsection{The Balance Condition}

At the maximum, moving slightly to the right should not improve the product.

Also, moving slightly to the left should not improve the product.

So at the best point, there should be no first-order advantage in either direction.

That means the main part of the change must vanish:

\[
L-2x=0.
\]

Solving,

\[
2x=L.
\]

Hence

\[
x=\frac L2.
\]

Therefore,

\[
AC=\frac L2
\]

and

\[
CB=\frac L2.
\]

So the product $(AC\cdot CB)$ is largest when (C) is the midpoint of (AB).

\subsection{What Was Discovered}

We did not begin with the derivative.

We began with a concrete problem.

Then we asked what happens when the input changes slightly.

This forced us to compare

\[
P(x+h)
\]

with

\[
P(x).
\]

The difference

\[
P(x+h)-P(x)
\]

measured the change in the quantity.

Then we found that the most important part of the change was the part proportional to (h):

\[
h(L-2x).
\]

At the maximum, this first-order part had to disappear.

This is the beginning of the derivative.

\begin{quote}
The derivative is born from the attempt to understand which part of a small change controls increase and decrease.
\end{quote}

\subsection{Challenge}

Try the same idea with

\[
P(x)=x^2.
\]

Compute

\[
P(x+h)-P(x).
\]

Which term is proportional to (h)?

Which term is proportional to $(h^2)$?

What expression controls the first-order change?

\subsection{Challenge}

Now try

\[
P(x)=x^3.
\]

Compute

\[
P(x+h)-P(x).
\]

Expand:

\[
(x+h)^3-x^3.
\]

Which part is proportional to (h)?

What expression appears multiplying (h)?

This expression will later become the derivative of $(x^3)$.


\section{The Need for a Derivative}

Continuity tells us when a function has no breaks.

But continuity does not tell us how fast the function changes.

Consider a moving object.

At time (x), suppose its position is

\[
y=f(x).
\]

If time changes from (x) to $(x+\Delta x)$, then the position changes from

\[
f(x)
\]

to

\[
f(x+\Delta x).
\]

So the change in position is

\[
\Delta y=f(x+\Delta x)-f(x).
\]

The average rate of change is

\[
\frac{\Delta y}{\Delta x} = \frac{f(x+\Delta x)-f(x)}{\Delta x}.
\]

This is the slope of the secant line between the two points

\[
(x,f(x))
\]

and

\[
(x+\Delta x,f(x+\Delta x)).
\]

But now we ask a sharper question.

\begin{quote}
What is the rate of change at one instant?
\end{quote}

An instant has no length, so we cannot simply divide by zero.

Instead, we take smaller and smaller intervals and look for the value approached by the average rates of change.

This creates the derivative.

\begin{definition}
Let (y=f(x)). For an increment $(\Delta x\neq 0)$, define

\[
\Delta y=f(x+\Delta x)-f(x).
\]

If the limit

\[
\lim_{\Delta x\to 0}
\frac{\Delta y}{\Delta x}
\]

exists, then this limit is called the \emph{derivative} of (f) at (x), and is denoted by

\[
f'(x)
\]

or

\[
\frac{dy}{dx}.
\]

Thus,

\[
f'(x) = \lim_{\Delta x\to 0} \frac{f(x+\Delta x)-f(x)}{\Delta x}.
\]
\end{definition}

\subsection{Example: The Derivative of $(x^2)$}

Let

\[
f(x)=x^2.
\]

Then

\[
f(x+\Delta x)=(x+\Delta x)^2.
\]

So

\[
\Delta y = f(x+\Delta x)-f(x) (x+\Delta x)^2-x^2.
\]

Expanding,

\[
(x+\Delta x)^2=x^2+2x\Delta x+(\Delta x)^2.
\]

Therefore,

\[
\Delta y = x^2+2x\Delta x+(\Delta x)^2-x^2 2x\Delta x+(\Delta x)^2.
\]

Now divide by $(\Delta x)$:

\[
\frac{\Delta y}{\Delta x} = \frac{2x\Delta x+(\Delta x)^2}{\Delta x} 2x+\Delta x.
\]

Taking the limit as $(\Delta x\to 0)$,

\[
f'(x) = \lim_{\Delta x\to 0}(2x+\Delta x)2x.
\]

Thus,

\[
\frac{d}{dx}(x^2)=2x.
\]

\subsection{What We Have Discovered}

We began with quantities that approach values without necessarily reaching them.

This forced us to define limits.

Then we noticed that many differences approach zero.

This forced us to name infinitesimals.

Then we asked when the approached value of a function equals the actual value.

This forced us to define continuity.

Finally, we asked how to measure instantaneous change.

This forced us to define the derivative.

\begin{quote}
The derivative is not a formula first. It is the answer to the problem of measuring change at an instant.
\end{quote}

\subsection{Problems}

\subsubsection*{Problem 8}

Using the definition of derivative, compute the derivative of

\[
f(x)=x.
\]

\subsubsection*{Problem 9}

Using the definition of derivative, compute the derivative of

\[
f(x)=x^3.
\]

\subsubsection*{Problem 10}

Let

\[
f(x)=|x|.
\]

Compute the average rate of change

\[
\frac{f(0+\Delta x)-f(0)}{\Delta x}.
\]

What happens when $(\Delta x>0)$?

What happens when $(\Delta x<0)$?

Does the derivative exist at (x=0)?
<<<<<<< HEAD


\section{A Problem Before the Integral: Measuring a Curved Area}

The derivative was born from a question:

\begin{quote}
What is the instantaneous rate of change?
\end{quote}

Now we ask an opposite question.

Suppose many small quantities are being added together.

\begin{quote}
How can many tiny pieces produce a total quantity?
\end{quote}

This question leads to the integral.

\subsection{The Area Problem}

Consider the function

\[
f(x)=x^2.
\]

We want to find the area under the curve from (x=0) to (x=1).

That is, we want the area of the region bounded by

\[
y=x^2,
\]

the (x)-axis, and the vertical lines

\[
x=0
\]

and

\[
x=1.
\]

For a rectangle, area is easy:

\[
\text{area}=\text{base}\cdot\text{height}.
\]

But this region is not a rectangle.

Its upper boundary is curved.

So we ask:

\begin{quote}
Can a curved region be approximated by many thin rectangles?
\end{quote}

\subsection{First Approximation}

Divide the interval ([0,1]) into (n) equal parts.

Each small interval has width

\[
\Delta x=\frac{1}{n}.
\]

The division points are

\[
0,\frac{1}{n},\frac{2}{n},\frac{3}{n},\ldots,\frac{n}{n}.
\]

On the (k)-th interval, use the right endpoint

\[
x_k=\frac{k}{n}.
\]

The height of the rectangle is then

\[
f(x_k)=\left(\frac{k}{n}\right)^2.
\]

So the area of the (k)-th rectangle is

\[
f(x_k)\Delta x = \left(\frac{k}{n}\right)^2\frac{1}{n}.
\]

Therefore, the total approximate area is

\[
A_n=
\sum_{k=1}^{n}
\left(\frac{k}{n}\right)^2\frac{1}{n}.
\]

Simplifying,

\[
A_n=
\frac{1}{n^3}
\sum_{k=1}^{n}k^2.
\]

Now use the known formula

\[
\sum_{k=1}^{n}k^2 = \frac{n(n+1)(2n+1)}{6}.
\]

Thus,

\[
A_n=
\frac{1}{n^3}\cdot
\frac{n(n+1)(2n+1)}{6}.
\]

So

\[
A_n=
\frac{(n+1)(2n+1)}{6n^2}.
\]

Expanding,

\[
A_n=
\frac{2n^2+3n+1}{6n^2}.
\]

Therefore,

\[
A_n=
\frac{1}{3}+\frac{1}{2n}+\frac{1}{6n^2}.
\]

As $(n\to\infty)$,

\[
\frac{1}{2n}\to 0
\]

and

\[
\frac{1}{6n^2}\to 0.
\]

Hence,

\[
\lim_{n\to\infty}A_n = \frac{1}{3}.
\]

So the area under $(y=x^2)$ from (0) to (1) is

\[
\frac{1}{3}.
\]

\subsection{What Was Discovered}

We did not begin with the integral.

We began with a problem:

\begin{quote}
How do we measure an area whose boundary is curved?
\end{quote}

A curved region was difficult, but rectangles were easy.

So we replaced the curved region by many rectangles.

Each rectangle had area

\[
f(x_k)\Delta x.
\]

Adding all rectangles gave

\[
\sum_{k=1}^{n} f(x_k)\Delta x.
\]

Then we made the rectangles thinner and thinner by taking

\[
n\to\infty.
\]

This suggests a new object.

\begin{definition}
Let (f(x)) be a function on ([a,b]). Divide the interval into (n) small parts of width

\[
\Delta x=\frac{b-a}{n}.
\]

Choose points $(x_k)$ in each small interval. If the limit

\[
\lim_{n\to\infty}
\sum_{k=1}^{n} f(x_k)\Delta x
\]

exists, then this limit is called the \emph{definite integral} of (f) from (a) to (b), and is written

\[
\int_a^b f(x),dx.
\]
\end{definition}

Thus, the symbol

\[
\int_a^b f(x),dx
\]

means:

\begin{quote}
the limit of a sum of infinitely many infinitesimal contributions.
\end{quote}

\subsection{The Integral of $(x^2)$}

From our computation,

\[
\int_0^1 x^2,dx = \frac{1}{3}.
\]

This was not magic.

It came from the limit

\[
\int_0^1 x^2,dx = \lim_{n\to\infty}
\sum_{k=1}^{n}
\left(\frac{k}{n}\right)^2\frac{1}{n}.
\]

So the integral is born from approximation.

\begin{quote}
The integral is not first a formula. It is the answer to the problem of measuring a total made from many tiny parts.
\end{quote}
=======
>>>>>>> eeca8e76b26e64e46bf02563453f92c0884596ae
